\section{DFA minimization and the Myhill--Nerode theorem}

\subsection{Equivalent versus distinguishable states}
Let $M=(Q,\Sigma,\transition,q_0,F)$ be a DFA. Write
$\transition^*(q,x)$ for the state reached from $q$ on $x$.

\begin{definition}[Distinguishable states]
States $p,q\in Q$ are \emph{equivalent} if
\[
  \forall x\in\Sigma^*:\quad
  \bigl(\transition^*(p,x)\in F
        \iff
        \transition^*(q,x)\in F\bigr).
\]
Otherwise they are \emph{distinguishable}: some witness $x$ sends exactly one
of them into $F$. Equivalently, once the machine is in $p$ versus $q$, no
future input can tell the two situations apart if and only if $p$ and $q$
are equivalent.
\end{definition}
Accepting and non-accepting states are always distinguishable (witness $\eps$).

\subsection{The table-filling algorithm}
\begin{algorithm}[Table filling]
\begin{enumerate}
  \item Remove unreachable states by BFS/DFS from $q_0$.
  \item Make a table of unordered pairs $\set{p,q}$ with $p\neq q$.
  \item Mark every pair in which exactly one state is accepting.
  \item Repeat: if an unmarked pair $\set{p,q}$ has, for some $a\in\Sigma$,
        the pair $\set{\transition(p,a),\transition(q,a)}$ already marked, then
        mark $\set{p,q}$. (If $\transition(p,a)=r$ and $\transition(q,a)=s$ are
        distinguishable, then $ax$ distinguishes $p$ and $q$ whenever $x$
        distinguishes $r$ and $s$.)
  \item Stop when a pass marks nothing new. Unmarked pairs are equivalent.
        Merge equivalent states.
\end{enumerate}
\end{algorithm}

\begin{theorem}[Correctness of table filling]
A pair is marked if and only if it is distinguishable.
\end{theorem}

\begin{proof}
($\Rightarrow$) Induction on the iteration in which the pair is marked. Pairs
marked in step~3 are distinguished by $\eps$. If $\set{p,q}$ is marked because
$\set{\transition(p,a),\transition(q,a)}$ was already marked, the inductive
hypothesis gives a witness $x$ for the successors, so $ax$ witnesses
$\set{p,q}$.

($\Leftarrow$) Induction on the length of a shortest distinguishing word.
If $|x|=0$, then $x=\eps$ and the pair is marked in step~3. If
$x=ax'$ with $|x|=k+1$, then $x'$ distinguishes
$r=\transition(p,a)$ and $s=\transition(q,a)$ and has length $k$, so
$\set{r,s}$ is already marked by the inductive hypothesis. Step~4 therefore
marks $\set{p,q}$.
\end{proof}

The relation ``unmarked'' is an equivalence relation on the (reachable) states.
The \emph{minimal DFA} has those equivalence classes as states:
\[
  M'=\bigl(Q/{\sim},\;\Sigma,\;\transition',\;[q_0],\;
           \set{[q]:q\in F}\bigr),
  \qquad
  \transition'([q],a)=[\transition(q,a)].
\]
$\transition'$ is well-defined: if $q\sim r$ but
$\transition(q,a)\not\sim\transition(r,a)$, then some $x$ distinguishes the
successors, so $ax$ distinguishes $q$ and $r$. An accept state and a reject
state never lie in the same class.

\begin{theorem}
$\lang{M'}=\lang{M}$, and $M'$ has the fewest states among all DFAs equivalent
to $M$.
\end{theorem}

\begin{proof}[Sketch]
A run of $M$ on $w$ visits $q_0,q_1,\ldots,q_n$; the corresponding run of $M'$
visits $[q_0],[q_1],\ldots,[q_n]$. The last class is accepting if and only if
$q_n\in F$, so the languages coincide.

If some equivalent DFA $N$ had fewer states than $M'$, then two distinct
classes $[p]$ and $[q]$ of $M'$ would map to the same state of $N$. A witness
that distinguishes $p$ and $q$ would then distinguish that state of $N$ from
itself, which is impossible. Hence $M'$ is minimal.
\end{proof}

\begin{example}[Table filling on a four-state DFA]
Let $Q=\set{1,2,3,4}$ with $1$ start, $F=\set{3,4}$, and
\[
\begin{array}{c|cc}
 & 0 & 1\\
\hline
1 & 2 & 3\\
2 & 1 & 4\\
3 & 4 & 3\\
4 & 3 & 4
\end{array}
\]
Pairs mixing $\set{1,2}$ with $\set{3,4}$ are marked immediately. Then
$\transition(1,1)=3$ and $\transition(2,1)=4$ are both accepting, while
$\transition(3,0)=4$ and $\transition(4,0)=3$ stay unmarked. The unmarked
pairs are $\set{1,2}$ and $\set{3,4}$. The minimal DFA therefore has two
states: $\set{1,2}$ (reject) and $\set{3,4}$ (accept).
\end{example}

\subsection{Equivalence of strings}
The same idea, lifted from states to strings, is independent of any particular
machine.

\begin{definition}[Myhill--Nerode relation]
For a language $A\subseteq\Sigma^*$ define
\[
  x\equiv_A y
  \quad\Longleftrightarrow\quad
  \forall z\in\Sigma^*:\;
  (xz\in A\iff yz\in A).
\]
Strings $x,y$ are equivalent when no future $z$ can tell them apart; they
represent the same knowledge about $A$.
\end{definition}

\begin{example}[Strings ending in $1$]
Let $L=\set{w:w\text{ ends in }1}$. Then
$00\equiv_L 000\equiv_L 1000$ (all ``last symbol is $0$''), while
$1$, $01$, $001$ are another class. Two classes: last bit $0$ versus last
bit $1$.

For $A=\set{w:w\text{ ends in }00}$ there are three relevant classes:
ends in $00$; ends in $1$ (or is $\eps$ / ends in a $1$ not covered); ends
in a single $0$ not already forming the suffix $00$.
\end{example}

\subsection{Myhill--Nerode theorem}
\begin{theorem}[Myhill--Nerode]
The following are equivalent for $A\subseteq\Sigma^*$.
\begin{enumerate}
  \item $A$ is regular.
  \item $\equiv_A$ has finitely many equivalence classes.
\end{enumerate}
Moreover, the number of classes equals the number of states in the unique
(up to renaming) minimal DFA for $A$.
\end{theorem}

\begin{proof}
($\Rightarrow$) Let $M=(Q,\Sigma,\transition,q_0,F)$ recognize $A$. If
$\transition^*(q_0,x)=\transition^*(q_0,y)$, then $x\equiv_A y$, because
every continuation $z$ leads both strings to the same state. Contrapositively,
distinguishable strings go to different states. Hence the number of classes is
at most $|Q|$, which is finite.

($\Leftarrow$) Let the classes of $\equiv_A$ be the states of a DFA
\[
  M=\bigl(\set{[x]:x\in\Sigma^*},\;
          \Sigma,\;
          \transition,\;
          [\eps],\;
          \set{[x]:x\in A}\bigr),
\]
where $\transition([x],a)=[xa]$. This is well-defined by right-invariance of
$\equiv_A$: if $x\equiv_A y$ then $xa\equiv_A ya$.

\textbf{Claim.} $\transition^*([\eps],w)=[w]$ for every $w$. By induction on
$|w|$: the base is $[\eps]$; if the claim holds for $w$ then
\[
  \transition^*([\eps],wa)
    =\transition\bigl(\transition^*([\eps],w),a\bigr)
    =\transition([w],a)=[wa].
\]
Hence $M$ accepts $w$ if and only if $[w]$ is accepting if and only if
$w\in A$.

The construction uses exactly one state per class, so it is minimal. It
matches the table-filling DFA: two reachable states $p=\transition^*(q_0,x)$
and $q=\transition^*(q_0,y)$ are equivalent if and only if $x\equiv_A y$.
\end{proof}
